% =====================================================================
% Capítulo de Resultados — Etapa 4 (integrado, Frentes A/B/C do Overleaf)
% ---------------------------------------------------------------------
% COMO USAR NO OVERLEAF:
%   1. Faça upload deste arquivo e das figuras (pasta figuras/) no projeto.
%   2. No preâmbulo, garanta: \usepackage{graphicx} e \usepackage{float}
%      (já presentes no seu documento) e, se as figuras ficarem em subpasta:
%         \graphicspath{{figuras/}}
%   3. Antes de \bibliographystyle, insira:  % =====================================================================
% Capítulo de Resultados — Etapa 4 (integrado, Frentes A/B/C do Overleaf)
% ---------------------------------------------------------------------
% COMO USAR NO OVERLEAF:
%   1. Faça upload deste arquivo e das figuras (pasta figuras/) no projeto.
%   2. No preâmbulo, garanta: \usepackage{graphicx} e \usepackage{float}
%      (já presentes no seu documento) e, se as figuras ficarem em subpasta:
%         \graphicspath{{figuras/}}
%   3. Antes de \bibliographystyle, insira:  % =====================================================================
% Capítulo de Resultados — Etapa 4 (integrado, Frentes A/B/C do Overleaf)
% ---------------------------------------------------------------------
% COMO USAR NO OVERLEAF:
%   1. Faça upload deste arquivo e das figuras (pasta figuras/) no projeto.
%   2. No preâmbulo, garanta: \usepackage{graphicx} e \usepackage{float}
%      (já presentes no seu documento) e, se as figuras ficarem em subpasta:
%         \graphicspath{{figuras/}}
%   3. Antes de \bibliographystyle, insira:  % =====================================================================
% Capítulo de Resultados — Etapa 4 (integrado, Frentes A/B/C do Overleaf)
% ---------------------------------------------------------------------
% COMO USAR NO OVERLEAF:
%   1. Faça upload deste arquivo e das figuras (pasta figuras/) no projeto.
%   2. No preâmbulo, garanta: \usepackage{graphicx} e \usepackage{float}
%      (já presentes no seu documento) e, se as figuras ficarem em subpasta:
%         \graphicspath{{figuras/}}
%   3. Antes de \bibliographystyle, insira:  \input{resultados}
%
% Este arquivo SUBSTITUI o antigo frente_b_engenharia_migracao.tex: a Frente B
% já está incluída aqui como subseção (5.3), com as mesmas tabelas reais.
% Chaves de citação usadas já existem no seu .bib: oliveira2023, rampinelli2026,
% nielsen1994, cardieri2018.
% =====================================================================

\section{Resultados Obtidos}
\label{sec:resultados}

Este capítulo apresenta os resultados da avaliação tripartite do artefato. Por
transparência metodológica, distinguem-se dois tipos de resultado: os
\textbf{já apurados} --- a Frente~B (engenharia da migração), derivada da análise
direta do código-fonte e, portanto, determinística e reprodutível --- e os
\textbf{pendentes de coleta empírica} --- a Frente~A (desempenho), que depende dos
ensaios com ADB em dispositivo físico, e a Frente~C (usabilidade), que depende das
inspeções dos especialistas. Para as frentes pendentes, apresentam-se o
instrumento de medição já construído, tabelas preparadas para os dados e uma
demonstração do método com dados sintéticos, explicitamente rotulados como
ilustrativos.

\subsection{O produto: aplicativo móvel}
\label{sub:produto}

O primeiro resultado é o próprio artefato: o cliente móvel em React Native/Expo,
que consome a mesma API FastAPI do PWA e preserva as funcionalidades da versão
web (entrada manual de valores, envio de laudo em PDF, verificação de
disponibilidade do servidor, exibição de resultados/\emph{briefing} e tabela de
comparação de referências). A caracterização técnica dessa migração constitui a
Frente~B (Seção~\ref{sub:frenteB}).

\subsection{Frente A --- Desempenho comparativo (ADB)}
\label{sub:frenteA}

A Frente~A compara PWA e cliente móvel sobre as mesmas tarefas (T1 inicialização;
T2 submissão de PDF; T3 rolagem; T4 renderização de tabelas), no mesmo
dispositivo Android, com dez execuções por tarefa e cliente, alternadas para
diluir efeitos térmicos e de processos concorrentes. As métricas --- tempo de
resposta, uso de CPU, memória (RAM), taxa de quadros (FPS) e energia estimada ---
são coletadas via Android Debug Bridge (ADB) e \texttt{dumpsys}, conforme o
protocolo de Oliveira et al. (2023) \cite{oliveira2023}.

Os ensaios foram conduzidos em um tablet Samsung Galaxy Tab~S9~FE
(modelo SM-X510), com Android~16, com dez repetições por medição.
A Tabela~\ref{tab:desempenho} consolida os resultados obtidos. O achado mais
expressivo diz respeito ao consumo de memória: o cliente móvel nativo utilizou,
em média, \textbf{84,9~MB}, ante \textbf{235,0~MB} do PWA executado no
navegador --- uma redução de \textbf{63,9\%}, isto é, o PWA consome cerca de
\textbf{2,8 vezes mais} memória para a mesma funcionalidade. A inicialização a
frio do aplicativo nativo situou-se em \textbf{241~ms} (desvio-padrão de
33~ms). A fluidez da rolagem foi aferida com o \texttt{dumpsys gfxinfo}: o cliente nativo
manteve a rolagem praticamente sem travamentos, com apenas \textbf{0,2\%} de
quadros com atraso (\emph{janky}) e tempo de renderização no 95º percentil em torno
de 9--12~ms, dentro do orçamento do \emph{display}. Esse resultado reproduziu-se em duas coletas independentes (0,2\% de quadros \emph{janky} em ambas; $n=20$ no total), reforçando sua robustez. A mesma medição não pôde ser
obtida para o PWA: o Chrome realiza a composição gráfica por um \emph{pipeline}
próprio, que não expõe essas estatísticas ao \texttt{gfxinfo} --- uma assimetria de
medição já prevista para a comparação PWA $\times$ nativo. O uso de CPU durante a
rolagem manteve-se baixo em ambos os clientes, porém o instrumento
(\texttt{dumpsys cpuinfo}) não produziu variação apreciável entre repetições, não
permitindo comparação conclusiva nesse ponto. O consumo energético estimado,
também previsto no protocolo, não foi aferido nesta etapa, permanecendo como
trabalho futuro.

\begin{table}[H]
\centering
\caption{Desempenho PWA $\times$ Móvel (Galaxy Tab~S9~FE, Android~16; $n=10$; média $\pm$ desvio-padrão).}
\label{tab:desempenho}
\footnotesize
\begin{tabular}{|p{6.2cm}|c|c|c|}
\hline
\textbf{Métrica (direção)} & \textbf{PWA} & \textbf{Móvel} & \textbf{$\Delta$} \\ \hline
Memória RAM --- app carregado (MB, $\downarrow$) & $235{,}0 \pm 18{,}7$ & $84{,}9 \pm 0{,}2$ & $-63{,}9\%$ \\ \hline
Inicialização a frio (ms, $\downarrow$) & --- & $241{,}3 \pm 33{,}3$ & --- \\ \hline
Fluidez em rolagem --- quadros \emph{janky} (\%, $\downarrow$) & n/d (ver texto) & $0{,}2$ & --- \\ \hline
CPU em rolagem (\%, $\downarrow$) & \multicolumn{3}{c|}{baixo em ambos; sem diferença conclusiva} \\ \hline
\end{tabular}
\end{table}

A Figura~\ref{fig:desempenho} apresenta os resultados por métrica (um painel por
métrica; barras = média, hastes = desvio-padrão). O contraste de memória entre o
PWA e o cliente nativo é expressivo e corrobora, com dados primários de um caso
real, a tese de que o paradigma \emph{WebView} impõe maior consumo de recursos
\cite{karin2025, sokolov2025}.

\begin{figure}[H]
\centering
\includegraphics[width=0.95\textwidth]{fig_frente_a_desempenho_adb.png}
\caption{Desempenho PWA $\times$ Móvel por métrica (Galaxy Tab~S9~FE, Android~16; $n=10$; barras = média, hastes = desvio-padrão).}
\label{fig:desempenho}
\end{figure}

\subsection{Frente B --- Engenharia da migração}
\label{sub:frenteB}

A Frente~B caracteriza o reuso arquitetural viabilizado pela estratégia
\emph{API-first}, a partir do mapeamento de migração das duas bases de código do
cliente (PWA em React/TypeScript e aplicativo em React Native/Expo) e do serviço
compartilhado (\emph{backend} FastAPI). Conforme o procedimento descrito na
Seção~4.3, cada artefato preexistente foi classificado segundo três padrões de
portabilidade: \textbf{(i)}~reescrita completa; \textbf{(ii)}~portabilidade
estrutural com adaptação de elementos visuais; e \textbf{(iii)}~substituição por
componente nativo equivalente. As métricas de linhas de código (LOC) foram
extraídas diretamente do repositório.

\subsubsection{Reuso do núcleo funcional via API-first}

O resultado central é o \textbf{reuso integral (100\%) de toda a lógica
clinicamente sensível}. O \emph{backend} FastAPI --- extração por OCR/\emph{parsing},
classificação pelas faixas da PNS, seleção de especialidades e geração do
\emph{briefing} --- foi reaproveitado \textbf{sem uma única linha reescrita}, sendo
consumido pelo novo cliente pelos mesmos três \emph{endpoints} REST
(Tabela~\ref{tab:endpoints}). A base de conhecimento (\texttt{data/*.csv}) também
permaneceu inalterada. Em números, \textbf{1.654 linhas de Python} e \textbf{145
linhas de regras em CSV} foram preservadas na íntegra \cite{rampinelli2026}.

\begin{table}[H]
\centering
\caption{\emph{Endpoints} REST reaproveitados sem alteração pelos dois clientes.}
\label{tab:endpoints}
\footnotesize
\begin{tabular}{|l|l|p{6.2cm}|}
\hline
\textbf{Endpoint} & \textbf{Método} & \textbf{Função (reusada por PWA e Mobile)} \\ \hline
\texttt{/health} & GET & Verificação de disponibilidade do serviço \\ \hline
\texttt{/interpret} & POST & Interpretação a partir de laudo em PDF (multipart) \\ \hline
\texttt{/interpret-manual} & POST & Interpretação a partir de valores digitados (JSON) \\ \hline
\end{tabular}
\end{table}

\subsubsection{Mapa de migração dos componentes}

A Tabela~\ref{tab:mapa_migracao} apresenta o mapa de migração, artefato a
artefato. A maioria dos componentes de interface enquadrou-se no padrão~(ii):
estrutura declarativa e lógica de estado preservadas, com adaptação apenas das
primitivas visuais (de \texttt{<div>}/\texttt{<input>} para
\texttt{<View>}/\texttt{<TextInput>}/\texttt{StyleSheet}). As substituições
nativas (iii) concentraram-se nos pontos de contato com \emph{hardware} e ambiente.

\begin{table}[H]
\centering
\caption{Mapa de migração PWA $\rightarrow$ Mobile por artefato (LOC = linhas de código).}
\label{tab:mapa_migracao}
\scriptsize
\setlength{\tabcolsep}{3pt}
\begin{tabular}{|p{4.3cm}|p{4.0cm}|c|p{3.0cm}|}
\hline
\textbf{Artefato na PWA (LOC)} & \textbf{Contraparte no Mobile} & \textbf{LOC} & \textbf{Padrão} \\ \hline
\emph{Backend} FastAPI + serviços (1.654) & idêntico (via REST) & 1.654 & Reuso integral \\ \hline
Base de conhecimento \texttt{data/*.csv} (145) & idêntica (via REST) & 145 & Reuso integral \\ \hline
\texttt{types/index.ts} (56) & \texttt{types.ts} (58) & 58 & (ii) Estrutural (contrato) \\ \hline
\texttt{services/api.ts} (213) & \texttt{api.ts} (71) & 71 & (ii) Estrutural c/ adaptação de transporte \\ \hline
\texttt{App.tsx} (323) & \texttt{App.tsx} (209) & 209 & (ii) Estrutural (orquestração/estado) \\ \hline
\texttt{ManualEntryForm.tsx} (90) & \texttt{ManualEntryForm.tsx} & --- & (ii) Estrutural \\ \hline
\texttt{ResultsDisplay.tsx} (196) & \texttt{Results.tsx} & --- & (ii) Estrutural \\ \hline
\texttt{ReferenceComparison.tsx} (82) & \texttt{ReferenceComparisonTable.tsx} & --- & (ii) Estrutural \\ \hline
\texttt{FileUpload.tsx} (56, \emph{react-dropzone}) & \texttt{PdfUpload.tsx} (\emph{expo-document-picker}) & --- & (iii) Substituição nativa \\ \hline
\texttt{LoadingSpinner.tsx} (79) & \texttt{ActivityIndicator} (nativo) & --- & (iii) Substituição nativa \\ \hline
\texttt{index.tsx} (66, \emph{ReactDOM}) & \texttt{index.ts} (8, Expo) & 8 & (iii) \emph{Bootstrap} de plataforma \\ \hline
\texttt{utils/} (1.239: contorno de \emph{webview}) & --- (não portados) & 0 & Eliminados \\ \hline
\end{tabular}
\end{table}

\subsubsection{Taxa de reuso e eliminação de código de contorno}

A Tabela~\ref{tab:taxa_reuso} sintetiza a taxa de reuso por camada. Destaca-se um
achado quantitativo relevante: a PWA mantinha três módulos utilitários
(\texttt{polyfills.ts}, \texttt{mobileOptimizations.ts} e
\texttt{mobileDetection.ts}), somando \textbf{1.239 linhas --- 49\% de todo o
código do cliente web (2.525 LOC)} ---, que existiam exclusivamente para mitigar
limitações do paradigma \emph{webview} em navegadores móveis. Na renderização
nativa do React Native, esses contornos tornaram-se desnecessários e foram
\textbf{integralmente eliminados}. Esse resultado é o correlato, no plano do
código, do \emph{overhead} de recursos que a literatura atribui às soluções
baseadas em \emph{webview} \cite{oliveira2023}.

\begin{table}[H]
\centering
\caption{Taxa de reuso por camada na migração PWA $\rightarrow$ Mobile.}
\label{tab:taxa_reuso}
\footnotesize
\begin{tabular}{|p{4.6cm}|c|p{5.2cm}|}
\hline
\textbf{Camada} & \textbf{Reuso} & \textbf{Observação} \\ \hline
Lógica de negócio (\emph{backend} + regras) & 100\% & 1.654 LOC Python + 145 CSV, sem alteração \\ \hline
Contrato de dados (interfaces TS) & $\approx$100\% & 5 interfaces de domínio migradas verbatim \\ \hline
Serviços de comunicação (rede) & 3/3 operações & Semântica preservada; transporte adaptado \\ \hline
Apresentação (componentes de UI) & Reescrita estrutural & Padrão (ii): estrutura/estado preservados \\ \hline
Utilitários de contorno de \emph{webview} & 0\% (eliminados) & 1.239 LOC descartadas \\ \hline
\end{tabular}
\end{table}

Confirma-se, assim, a hipótese de reuso: a separação de responsabilidades
\emph{API-first} preservou o núcleo funcional (ativo de maior valor e risco) e
localizou o esforço de reengenharia na camada de apresentação
\cite{rampinelli2026, oliveira2023}.

\subsection{Frente C --- Usabilidade (inspeção heurística)}
\label{sub:frenteC}

A Frente~C avaliou a usabilidade por inspeção heurística \cite{nielsen1994,
cardieri2018}: avaliadores inspecionaram, de forma independente, os dois clientes
(PWA e aplicativo) sobre os mesmos cenários, registrando violações às dez
heurísticas e classificando cada problema em uma escala de severidade de 0 a 4, no
modelo de \emph{um registro por problema}. Em razão do número de avaliadores
efetivamente reunido, esta frente é reportada como \textbf{estudo-piloto} de
caráter exploratório: participaram \textbf{quatro avaliadores}, dos quais três
registraram problemas de usabilidade (severidade $\geq 1$) e um registrou uma
observação classificada como severidade~0 (não contabilizada como problema);
registros de teste foram descartados.

A Tabela~\ref{tab:usabilidade} sintetiza os resultados, a
Tabela~\ref{tab:inventario_c} detalha os problemas e a
Figura~\ref{fig:usabilidade} mostra a distribuição por heurística.

\begin{table}[H]
\centering
\caption{Síntese comparativa da inspeção heurística (estudo-piloto; $n=3$ problemas).}
\label{tab:usabilidade}
\footnotesize
\begin{tabular}{|p{6.8cm}|c|c|}
\hline
\textbf{Métrica} & \textbf{PWA} & \textbf{Móvel} \\ \hline
Problemas registrados (severidade $\geq 1$) & 2 & 1 \\ \hline
Índice médio de severidade & 1,5 & 2,0 \\ \hline
Problemas de severidade 3--4 & 0 & 0 \\ \hline
Heurísticas violadas & H7, H8 & H7 \\ \hline
\end{tabular}
\end{table}

\begin{table}[H]
\centering
\caption{Inventário dos problemas de usabilidade identificados (estudo-piloto).}
\label{tab:inventario_c}
\scriptsize
\setlength{\tabcolsep}{3pt}
\begin{tabular}{|c|c|c|c|p{4.9cm}|p{4.0cm}|}
\hline
\textbf{\#} & \textbf{Cliente} & \textbf{Heur.} & \textbf{Sev.} & \textbf{Problema (resumo)} & \textbf{Recomendação} \\ \hline
1 & PWA & H8 & 1 & Tela de resultados poluída: \emph{tag} ``Normal'' redundante em todos os \emph{cards}; faltam unidades e valores de referência; dados de infraestrutura visíveis; sem divisão por séries. & Limpar os \emph{cards}; incluir unidades e referências; agrupar por séries. \\ \hline
2 & PWA & H7 & 2 & Achado que diverge apenas da referência clássica (não da PNS) não é explicado ao usuário. & Explicar o achado também nesse caso de divergência. \\ \hline
3 & Móvel & H7 & 2 & Após anexar o PDF, a aba de digitação manual permanece habilitada, permitindo entrada redundante. & Desabilitar a digitação manual quando há PDF anexado. \\ \hline
\end{tabular}
\end{table}

\begin{figure}[H]
\centering
\includegraphics[width=0.9\textwidth]{fig_frente_c_inspecao_heuristica.png}
\caption{Frente~C --- problemas por heurística, PWA $\times$ Móvel (estudo-piloto; $n=3$ problemas).}
\label{fig:usabilidade}
\end{figure}

Não se identificou nenhum problema de severidade alta (3--4): as ocorrências
situaram-se entre cosmética (1) e menor (2), sugerindo ausência de barreiras
graves de usabilidade nos fluxos inspecionados. Os achados concentraram-se na
\textbf{tela de apresentação de resultados} e em aspectos de \textbf{fluxo}:
poluição visual e ausência de informações essenciais de leitura, como unidades e
valores de referência --- heurística de estética e design minimalista (H8) ---, e
comportamentos pouco flexíveis, como a falta de explicação para achados que
divergem apenas da referência clássica e a permanência da digitação manual após o
anexo do PDF --- heurística de flexibilidade e eficiência de uso (H7). Registrou-se
ainda, no aplicativo, uma observação de severidade~0 (a tela de resultados abrir no
rodapé em vez do topo), não contabilizada como problema. Embora o índice médio de
severidade tenha sido ligeiramente maior no aplicativo (2,0 contra 1,5 do PWA), o
número reduzido de registros não autoriza comparação quantitativa conclusiva entre
os clientes; no formato piloto, o valor desta frente reside na natureza
\textbf{qualitativa e acionável} dos achados, que apontam melhorias concretas de
interface e pavimentam uma inspeção heurística mais ampla como trabalho futuro.

\subsection{Síntese: resultados obtidos $\times$ hipóteses}
\label{sub:sintese}

A Tabela~\ref{tab:sintese} relaciona os resultados às hipóteses do trabalho. A
hipótese de \textbf{reuso arquitetural} (Frente~B) está empiricamente confirmada;
as hipóteses de \textbf{desempenho} (A) e \textbf{usabilidade} (C) têm instrumentos
validados e aguardam a coleta.

\begin{table}[H]
\centering
\caption{Resultados obtidos frente às hipóteses de pesquisa.}
\label{tab:sintese}
\footnotesize
\begin{tabular}{|p{5.6cm}|p{3.0cm}|p{5.0cm}|}
\hline
\textbf{Dimensão (hipótese)} & \textbf{Situação} & \textbf{Evidência} \\ \hline
Artefato móvel funcional & Alcançado & Seção~\ref{sub:produto} \\ \hline
Ganho de desempenho (A) & Confirmado (parcial) & Seção~\ref{sub:frenteA}: RAM $-$63,9\%; rolagem nativa fluida (0,2\% \emph{janky}) \\ \hline
Reuso arquitetural API-first (B) & Alcançado (real) & Seção~\ref{sub:frenteB}: 100\% da lógica; $-$1.239 LOC de \emph{webview} \\ \hline
Ganho de usabilidade (C) & Explorado (estudo-piloto) & Seção~\ref{sub:frenteC}: 3 problemas (sev. 1--2); achados qualitativos \\ \hline
\end{tabular}
\end{table}

\subsection{Discussão: confronto com a literatura}
\label{sub:discussao}

Os resultados da Frente~B dialogam diretamente com os trabalhos relacionados
(Seção~3) e permitem confrontar o observado com o esperado. De início, confirma-se
a expectativa, derivada da literatura de arquitetura \emph{API-first}
\cite{rampinelli2026}, de que a centralização da lógica no \emph{backend}
minimizaria o custo de uma migração de paradigma: a hipótese de reuso não apenas se
sustentou, como se mostrou integral (100\%) no núcleo funcional.

Esse achado oferece uma contrapartida quantitativa ao estudo qualitativo de Zhou
(2024). Enquanto os vinte desenvolvedores entrevistados por Zhou relataram a
migração entre plataformas e a manutenção como desafios recorrentes do
desenvolvimento \emph{cross-platform}, sem, contudo, quantificar o esforço nem
documentar um caso real de transição \cite{zhou2024}, a Frente~B evidencia que esse
atrito é fortemente modulado pela decisão arquitetural. No InterpreteLabBR, a
arquitetura \emph{API-first} converteu o que a literatura descreve como
``dificuldade de migração'' em um esforço mensurável e localizado --- 100\% de
reuso da lógica de negócio, sem reescrita, e a eliminação de 1.239 linhas de
código de contorno ---, deslocando o custo da transição da lógica crítica para a
camada de apresentação. Em outras palavras, sob \emph{API-first}, a migração
deixou de ameaçar o núcleo funcional e passou a ser, predominantemente, um
exercício de portabilidade de interface.

O mesmo resultado fornece uma evidência, no plano do código-fonte, para a tese de
Karin (2025). Em revisão sistemática, Karin posiciona as soluções baseadas em
\emph{WebView} --- incluindo PWAs --- como as de maior \emph{overhead}, atribuindo
o custo à camada intermediária de DOM e do navegador \cite{karin2025}. As 1.239
linhas de utilitários eliminadas na migração (\emph{polyfills}, detecção de
dispositivo móvel e otimizações de rolagem) existiam exclusivamente para compensar
limitações do \emph{WebView} em navegadores móveis e correspondiam a 49\% do código
do cliente web. Sua supressão integral na arquitetura de renderização nativa
materializa, no artefato, o \emph{overhead} que Karin descreve de forma agregada:
além do custo de execução (energia e CPU), o paradigma \emph{WebView} impunha um
\emph{custo de engenharia} --- quase metade da base de código do cliente dedicada a
contornos ---, aqui suprimido. O achado, assim, corrobora empiricamente, e em um
caso real de saúde pública, a conclusão da revisão de Karin.

Por fim, os achados de Sokolov (2025) orientam a leitura dos resultados da Frente~A. Em \emph{benchmark} controlado via ADB, Sokolov observou que a
implementação nativa superou o React Native em eficiência de CPU e, sobretudo, no
crescimento de memória ao longo do tempo em rotinas contínuas, com vantagens
pontuais do React Native em cenários específicos \cite{sokolov2025}. Espera-se,
portanto, que a vantagem do cliente nativo sobre o PWA seja mais nítida em métricas
sensíveis à execução sustentada --- notadamente a memória RAM sob interação
contínua (rolagem) --- e mais modesta em tarefas dominadas por I/O e rede, como a
submissão do PDF, cuja latência é governada pela API. Diferentemente de Sokolov,
cujo \emph{benchmark} é sintético (criptografia AES) e desvinculado de domínio, a
Frente~A avaliará essas métricas em um fluxo real de uso clínico; a confirmação ou
refutação desse padrão de consumo de RAM constituirá o contraste central entre o
esperado pela literatura e o obtido em campo. As primeiras medições já apontam
nessa direção: em tablet Android, o cliente nativo consumiu \textbf{63,9\% menos
memória} que o PWA (84,9~MB contra 235,0~MB; $n=10$), convergindo com a vantagem
de memória relatada por Sokolov; a rolagem do cliente nativo mostrou-se fluida (apenas 0,2\% de quadros
\emph{janky}), ao passo que a aferição equivalente no PWA esbarrou em uma limitação
do instrumento, já que o Chrome não expõe essas estatísticas ao \texttt{gfxinfo}.

\subsection{Validade, limitações e extensão}
\label{sub:validade}

A Frente~B baseia-se na inspeção direta do código-fonte, não sujeita a viés
amostral. As Frentes~A e C controlam variáveis por desenho (dispositivo único e
execução alternada em A; heurísticas consagradas e severidade padronizada em C).
Como limitações: (i)~a Frente~A foi medida em um único dispositivo --- um tablet
de gama intermediária (Galaxy Tab~S9~FE), que não representa o aparelho de entrada
típico do SUS ---, de modo que as métricas \emph{absolutas} não se generalizam a
outros perfis de \emph{hardware}, preservando-se, contudo, a validade do resultado
\emph{relativo} (PWA $\times$ nativo medidos no mesmo aparelho); (ii)~a Frente~B é uma reconstrução documental
retrospectiva do mapa de migração; e (iii)~a Frente~C capta problemas de
conformidade heurística por especialistas, não a aceitação subjetiva de usuários
finais, e foi conduzida como estudo-piloto com painel reduzido, o que limita a generalização dos achados. As conclusões relativas sobre
reuso e desempenho têm potencial de transferência para migrações arquiteturais
correlatas na \emph{stack} adotada.

%
% Este arquivo SUBSTITUI o antigo frente_b_engenharia_migracao.tex: a Frente B
% já está incluída aqui como subseção (5.3), com as mesmas tabelas reais.
% Chaves de citação usadas já existem no seu .bib: oliveira2023, rampinelli2026,
% nielsen1994, cardieri2018.
% =====================================================================

\section{Resultados Obtidos}
\label{sec:resultados}

Este capítulo apresenta os resultados da avaliação tripartite do artefato. Por
transparência metodológica, distinguem-se dois tipos de resultado: os
\textbf{já apurados} --- a Frente~B (engenharia da migração), derivada da análise
direta do código-fonte e, portanto, determinística e reprodutível --- e os
\textbf{pendentes de coleta empírica} --- a Frente~A (desempenho), que depende dos
ensaios com ADB em dispositivo físico, e a Frente~C (usabilidade), que depende das
inspeções dos especialistas. Para as frentes pendentes, apresentam-se o
instrumento de medição já construído, tabelas preparadas para os dados e uma
demonstração do método com dados sintéticos, explicitamente rotulados como
ilustrativos.

\subsection{O produto: aplicativo móvel}
\label{sub:produto}

O primeiro resultado é o próprio artefato: o cliente móvel em React Native/Expo,
que consome a mesma API FastAPI do PWA e preserva as funcionalidades da versão
web (entrada manual de valores, envio de laudo em PDF, verificação de
disponibilidade do servidor, exibição de resultados/\emph{briefing} e tabela de
comparação de referências). A caracterização técnica dessa migração constitui a
Frente~B (Seção~\ref{sub:frenteB}).

\subsection{Frente A --- Desempenho comparativo (ADB)}
\label{sub:frenteA}

A Frente~A compara PWA e cliente móvel sobre as mesmas tarefas (T1 inicialização;
T2 submissão de PDF; T3 rolagem; T4 renderização de tabelas), no mesmo
dispositivo Android, com dez execuções por tarefa e cliente, alternadas para
diluir efeitos térmicos e de processos concorrentes. As métricas --- tempo de
resposta, uso de CPU, memória (RAM), taxa de quadros (FPS) e energia estimada ---
são coletadas via Android Debug Bridge (ADB) e \texttt{dumpsys}, conforme o
protocolo de Oliveira et al. (2023) \cite{oliveira2023}.

Os ensaios foram conduzidos em um tablet Samsung Galaxy Tab~S9~FE
(modelo SM-X510), com Android~16, com dez repetições por medição.
A Tabela~\ref{tab:desempenho} consolida os resultados obtidos. O achado mais
expressivo diz respeito ao consumo de memória: o cliente móvel nativo utilizou,
em média, \textbf{84,9~MB}, ante \textbf{235,0~MB} do PWA executado no
navegador --- uma redução de \textbf{63,9\%}, isto é, o PWA consome cerca de
\textbf{2,8 vezes mais} memória para a mesma funcionalidade. A inicialização a
frio do aplicativo nativo situou-se em \textbf{241~ms} (desvio-padrão de
33~ms). A fluidez da rolagem foi aferida com o \texttt{dumpsys gfxinfo}: o cliente nativo
manteve a rolagem praticamente sem travamentos, com apenas \textbf{0,2\%} de
quadros com atraso (\emph{janky}) e tempo de renderização no 95º percentil em torno
de 9--12~ms, dentro do orçamento do \emph{display}. Esse resultado reproduziu-se em duas coletas independentes (0,2\% de quadros \emph{janky} em ambas; $n=20$ no total), reforçando sua robustez. A mesma medição não pôde ser
obtida para o PWA: o Chrome realiza a composição gráfica por um \emph{pipeline}
próprio, que não expõe essas estatísticas ao \texttt{gfxinfo} --- uma assimetria de
medição já prevista para a comparação PWA $\times$ nativo. O uso de CPU durante a
rolagem manteve-se baixo em ambos os clientes, porém o instrumento
(\texttt{dumpsys cpuinfo}) não produziu variação apreciável entre repetições, não
permitindo comparação conclusiva nesse ponto. O consumo energético estimado,
também previsto no protocolo, não foi aferido nesta etapa, permanecendo como
trabalho futuro.

\begin{table}[H]
\centering
\caption{Desempenho PWA $\times$ Móvel (Galaxy Tab~S9~FE, Android~16; $n=10$; média $\pm$ desvio-padrão).}
\label{tab:desempenho}
\footnotesize
\begin{tabular}{|p{6.2cm}|c|c|c|}
\hline
\textbf{Métrica (direção)} & \textbf{PWA} & \textbf{Móvel} & \textbf{$\Delta$} \\ \hline
Memória RAM --- app carregado (MB, $\downarrow$) & $235{,}0 \pm 18{,}7$ & $84{,}9 \pm 0{,}2$ & $-63{,}9\%$ \\ \hline
Inicialização a frio (ms, $\downarrow$) & --- & $241{,}3 \pm 33{,}3$ & --- \\ \hline
Fluidez em rolagem --- quadros \emph{janky} (\%, $\downarrow$) & n/d (ver texto) & $0{,}2$ & --- \\ \hline
CPU em rolagem (\%, $\downarrow$) & \multicolumn{3}{c|}{baixo em ambos; sem diferença conclusiva} \\ \hline
\end{tabular}
\end{table}

A Figura~\ref{fig:desempenho} apresenta os resultados por métrica (um painel por
métrica; barras = média, hastes = desvio-padrão). O contraste de memória entre o
PWA e o cliente nativo é expressivo e corrobora, com dados primários de um caso
real, a tese de que o paradigma \emph{WebView} impõe maior consumo de recursos
\cite{karin2025, sokolov2025}.

\begin{figure}[H]
\centering
\includegraphics[width=0.95\textwidth]{fig_frente_a_desempenho_adb.png}
\caption{Desempenho PWA $\times$ Móvel por métrica (Galaxy Tab~S9~FE, Android~16; $n=10$; barras = média, hastes = desvio-padrão).}
\label{fig:desempenho}
\end{figure}

\subsection{Frente B --- Engenharia da migração}
\label{sub:frenteB}

A Frente~B caracteriza o reuso arquitetural viabilizado pela estratégia
\emph{API-first}, a partir do mapeamento de migração das duas bases de código do
cliente (PWA em React/TypeScript e aplicativo em React Native/Expo) e do serviço
compartilhado (\emph{backend} FastAPI). Conforme o procedimento descrito na
Seção~4.3, cada artefato preexistente foi classificado segundo três padrões de
portabilidade: \textbf{(i)}~reescrita completa; \textbf{(ii)}~portabilidade
estrutural com adaptação de elementos visuais; e \textbf{(iii)}~substituição por
componente nativo equivalente. As métricas de linhas de código (LOC) foram
extraídas diretamente do repositório.

\subsubsection{Reuso do núcleo funcional via API-first}

O resultado central é o \textbf{reuso integral (100\%) de toda a lógica
clinicamente sensível}. O \emph{backend} FastAPI --- extração por OCR/\emph{parsing},
classificação pelas faixas da PNS, seleção de especialidades e geração do
\emph{briefing} --- foi reaproveitado \textbf{sem uma única linha reescrita}, sendo
consumido pelo novo cliente pelos mesmos três \emph{endpoints} REST
(Tabela~\ref{tab:endpoints}). A base de conhecimento (\texttt{data/*.csv}) também
permaneceu inalterada. Em números, \textbf{1.654 linhas de Python} e \textbf{145
linhas de regras em CSV} foram preservadas na íntegra \cite{rampinelli2026}.

\begin{table}[H]
\centering
\caption{\emph{Endpoints} REST reaproveitados sem alteração pelos dois clientes.}
\label{tab:endpoints}
\footnotesize
\begin{tabular}{|l|l|p{6.2cm}|}
\hline
\textbf{Endpoint} & \textbf{Método} & \textbf{Função (reusada por PWA e Mobile)} \\ \hline
\texttt{/health} & GET & Verificação de disponibilidade do serviço \\ \hline
\texttt{/interpret} & POST & Interpretação a partir de laudo em PDF (multipart) \\ \hline
\texttt{/interpret-manual} & POST & Interpretação a partir de valores digitados (JSON) \\ \hline
\end{tabular}
\end{table}

\subsubsection{Mapa de migração dos componentes}

A Tabela~\ref{tab:mapa_migracao} apresenta o mapa de migração, artefato a
artefato. A maioria dos componentes de interface enquadrou-se no padrão~(ii):
estrutura declarativa e lógica de estado preservadas, com adaptação apenas das
primitivas visuais (de \texttt{<div>}/\texttt{<input>} para
\texttt{<View>}/\texttt{<TextInput>}/\texttt{StyleSheet}). As substituições
nativas (iii) concentraram-se nos pontos de contato com \emph{hardware} e ambiente.

\begin{table}[H]
\centering
\caption{Mapa de migração PWA $\rightarrow$ Mobile por artefato (LOC = linhas de código).}
\label{tab:mapa_migracao}
\scriptsize
\setlength{\tabcolsep}{3pt}
\begin{tabular}{|p{4.3cm}|p{4.0cm}|c|p{3.0cm}|}
\hline
\textbf{Artefato na PWA (LOC)} & \textbf{Contraparte no Mobile} & \textbf{LOC} & \textbf{Padrão} \\ \hline
\emph{Backend} FastAPI + serviços (1.654) & idêntico (via REST) & 1.654 & Reuso integral \\ \hline
Base de conhecimento \texttt{data/*.csv} (145) & idêntica (via REST) & 145 & Reuso integral \\ \hline
\texttt{types/index.ts} (56) & \texttt{types.ts} (58) & 58 & (ii) Estrutural (contrato) \\ \hline
\texttt{services/api.ts} (213) & \texttt{api.ts} (71) & 71 & (ii) Estrutural c/ adaptação de transporte \\ \hline
\texttt{App.tsx} (323) & \texttt{App.tsx} (209) & 209 & (ii) Estrutural (orquestração/estado) \\ \hline
\texttt{ManualEntryForm.tsx} (90) & \texttt{ManualEntryForm.tsx} & --- & (ii) Estrutural \\ \hline
\texttt{ResultsDisplay.tsx} (196) & \texttt{Results.tsx} & --- & (ii) Estrutural \\ \hline
\texttt{ReferenceComparison.tsx} (82) & \texttt{ReferenceComparisonTable.tsx} & --- & (ii) Estrutural \\ \hline
\texttt{FileUpload.tsx} (56, \emph{react-dropzone}) & \texttt{PdfUpload.tsx} (\emph{expo-document-picker}) & --- & (iii) Substituição nativa \\ \hline
\texttt{LoadingSpinner.tsx} (79) & \texttt{ActivityIndicator} (nativo) & --- & (iii) Substituição nativa \\ \hline
\texttt{index.tsx} (66, \emph{ReactDOM}) & \texttt{index.ts} (8, Expo) & 8 & (iii) \emph{Bootstrap} de plataforma \\ \hline
\texttt{utils/} (1.239: contorno de \emph{webview}) & --- (não portados) & 0 & Eliminados \\ \hline
\end{tabular}
\end{table}

\subsubsection{Taxa de reuso e eliminação de código de contorno}

A Tabela~\ref{tab:taxa_reuso} sintetiza a taxa de reuso por camada. Destaca-se um
achado quantitativo relevante: a PWA mantinha três módulos utilitários
(\texttt{polyfills.ts}, \texttt{mobileOptimizations.ts} e
\texttt{mobileDetection.ts}), somando \textbf{1.239 linhas --- 49\% de todo o
código do cliente web (2.525 LOC)} ---, que existiam exclusivamente para mitigar
limitações do paradigma \emph{webview} em navegadores móveis. Na renderização
nativa do React Native, esses contornos tornaram-se desnecessários e foram
\textbf{integralmente eliminados}. Esse resultado é o correlato, no plano do
código, do \emph{overhead} de recursos que a literatura atribui às soluções
baseadas em \emph{webview} \cite{oliveira2023}.

\begin{table}[H]
\centering
\caption{Taxa de reuso por camada na migração PWA $\rightarrow$ Mobile.}
\label{tab:taxa_reuso}
\footnotesize
\begin{tabular}{|p{4.6cm}|c|p{5.2cm}|}
\hline
\textbf{Camada} & \textbf{Reuso} & \textbf{Observação} \\ \hline
Lógica de negócio (\emph{backend} + regras) & 100\% & 1.654 LOC Python + 145 CSV, sem alteração \\ \hline
Contrato de dados (interfaces TS) & $\approx$100\% & 5 interfaces de domínio migradas verbatim \\ \hline
Serviços de comunicação (rede) & 3/3 operações & Semântica preservada; transporte adaptado \\ \hline
Apresentação (componentes de UI) & Reescrita estrutural & Padrão (ii): estrutura/estado preservados \\ \hline
Utilitários de contorno de \emph{webview} & 0\% (eliminados) & 1.239 LOC descartadas \\ \hline
\end{tabular}
\end{table}

Confirma-se, assim, a hipótese de reuso: a separação de responsabilidades
\emph{API-first} preservou o núcleo funcional (ativo de maior valor e risco) e
localizou o esforço de reengenharia na camada de apresentação
\cite{rampinelli2026, oliveira2023}.

\subsection{Frente C --- Usabilidade (inspeção heurística)}
\label{sub:frenteC}

A Frente~C avaliou a usabilidade por inspeção heurística \cite{nielsen1994,
cardieri2018}: avaliadores inspecionaram, de forma independente, os dois clientes
(PWA e aplicativo) sobre os mesmos cenários, registrando violações às dez
heurísticas e classificando cada problema em uma escala de severidade de 0 a 4, no
modelo de \emph{um registro por problema}. Em razão do número de avaliadores
efetivamente reunido, esta frente é reportada como \textbf{estudo-piloto} de
caráter exploratório: participaram \textbf{quatro avaliadores}, dos quais três
registraram problemas de usabilidade (severidade $\geq 1$) e um registrou uma
observação classificada como severidade~0 (não contabilizada como problema);
registros de teste foram descartados.

A Tabela~\ref{tab:usabilidade} sintetiza os resultados, a
Tabela~\ref{tab:inventario_c} detalha os problemas e a
Figura~\ref{fig:usabilidade} mostra a distribuição por heurística.

\begin{table}[H]
\centering
\caption{Síntese comparativa da inspeção heurística (estudo-piloto; $n=3$ problemas).}
\label{tab:usabilidade}
\footnotesize
\begin{tabular}{|p{6.8cm}|c|c|}
\hline
\textbf{Métrica} & \textbf{PWA} & \textbf{Móvel} \\ \hline
Problemas registrados (severidade $\geq 1$) & 2 & 1 \\ \hline
Índice médio de severidade & 1,5 & 2,0 \\ \hline
Problemas de severidade 3--4 & 0 & 0 \\ \hline
Heurísticas violadas & H7, H8 & H7 \\ \hline
\end{tabular}
\end{table}

\begin{table}[H]
\centering
\caption{Inventário dos problemas de usabilidade identificados (estudo-piloto).}
\label{tab:inventario_c}
\scriptsize
\setlength{\tabcolsep}{3pt}
\begin{tabular}{|c|c|c|c|p{4.9cm}|p{4.0cm}|}
\hline
\textbf{\#} & \textbf{Cliente} & \textbf{Heur.} & \textbf{Sev.} & \textbf{Problema (resumo)} & \textbf{Recomendação} \\ \hline
1 & PWA & H8 & 1 & Tela de resultados poluída: \emph{tag} ``Normal'' redundante em todos os \emph{cards}; faltam unidades e valores de referência; dados de infraestrutura visíveis; sem divisão por séries. & Limpar os \emph{cards}; incluir unidades e referências; agrupar por séries. \\ \hline
2 & PWA & H7 & 2 & Achado que diverge apenas da referência clássica (não da PNS) não é explicado ao usuário. & Explicar o achado também nesse caso de divergência. \\ \hline
3 & Móvel & H7 & 2 & Após anexar o PDF, a aba de digitação manual permanece habilitada, permitindo entrada redundante. & Desabilitar a digitação manual quando há PDF anexado. \\ \hline
\end{tabular}
\end{table}

\begin{figure}[H]
\centering
\includegraphics[width=0.9\textwidth]{fig_frente_c_inspecao_heuristica.png}
\caption{Frente~C --- problemas por heurística, PWA $\times$ Móvel (estudo-piloto; $n=3$ problemas).}
\label{fig:usabilidade}
\end{figure}

Não se identificou nenhum problema de severidade alta (3--4): as ocorrências
situaram-se entre cosmética (1) e menor (2), sugerindo ausência de barreiras
graves de usabilidade nos fluxos inspecionados. Os achados concentraram-se na
\textbf{tela de apresentação de resultados} e em aspectos de \textbf{fluxo}:
poluição visual e ausência de informações essenciais de leitura, como unidades e
valores de referência --- heurística de estética e design minimalista (H8) ---, e
comportamentos pouco flexíveis, como a falta de explicação para achados que
divergem apenas da referência clássica e a permanência da digitação manual após o
anexo do PDF --- heurística de flexibilidade e eficiência de uso (H7). Registrou-se
ainda, no aplicativo, uma observação de severidade~0 (a tela de resultados abrir no
rodapé em vez do topo), não contabilizada como problema. Embora o índice médio de
severidade tenha sido ligeiramente maior no aplicativo (2,0 contra 1,5 do PWA), o
número reduzido de registros não autoriza comparação quantitativa conclusiva entre
os clientes; no formato piloto, o valor desta frente reside na natureza
\textbf{qualitativa e acionável} dos achados, que apontam melhorias concretas de
interface e pavimentam uma inspeção heurística mais ampla como trabalho futuro.

\subsection{Síntese: resultados obtidos $\times$ hipóteses}
\label{sub:sintese}

A Tabela~\ref{tab:sintese} relaciona os resultados às hipóteses do trabalho. A
hipótese de \textbf{reuso arquitetural} (Frente~B) está empiricamente confirmada;
as hipóteses de \textbf{desempenho} (A) e \textbf{usabilidade} (C) têm instrumentos
validados e aguardam a coleta.

\begin{table}[H]
\centering
\caption{Resultados obtidos frente às hipóteses de pesquisa.}
\label{tab:sintese}
\footnotesize
\begin{tabular}{|p{5.6cm}|p{3.0cm}|p{5.0cm}|}
\hline
\textbf{Dimensão (hipótese)} & \textbf{Situação} & \textbf{Evidência} \\ \hline
Artefato móvel funcional & Alcançado & Seção~\ref{sub:produto} \\ \hline
Ganho de desempenho (A) & Confirmado (parcial) & Seção~\ref{sub:frenteA}: RAM $-$63,9\%; rolagem nativa fluida (0,2\% \emph{janky}) \\ \hline
Reuso arquitetural API-first (B) & Alcançado (real) & Seção~\ref{sub:frenteB}: 100\% da lógica; $-$1.239 LOC de \emph{webview} \\ \hline
Ganho de usabilidade (C) & Explorado (estudo-piloto) & Seção~\ref{sub:frenteC}: 3 problemas (sev. 1--2); achados qualitativos \\ \hline
\end{tabular}
\end{table}

\subsection{Discussão: confronto com a literatura}
\label{sub:discussao}

Os resultados da Frente~B dialogam diretamente com os trabalhos relacionados
(Seção~3) e permitem confrontar o observado com o esperado. De início, confirma-se
a expectativa, derivada da literatura de arquitetura \emph{API-first}
\cite{rampinelli2026}, de que a centralização da lógica no \emph{backend}
minimizaria o custo de uma migração de paradigma: a hipótese de reuso não apenas se
sustentou, como se mostrou integral (100\%) no núcleo funcional.

Esse achado oferece uma contrapartida quantitativa ao estudo qualitativo de Zhou
(2024). Enquanto os vinte desenvolvedores entrevistados por Zhou relataram a
migração entre plataformas e a manutenção como desafios recorrentes do
desenvolvimento \emph{cross-platform}, sem, contudo, quantificar o esforço nem
documentar um caso real de transição \cite{zhou2024}, a Frente~B evidencia que esse
atrito é fortemente modulado pela decisão arquitetural. No InterpreteLabBR, a
arquitetura \emph{API-first} converteu o que a literatura descreve como
``dificuldade de migração'' em um esforço mensurável e localizado --- 100\% de
reuso da lógica de negócio, sem reescrita, e a eliminação de 1.239 linhas de
código de contorno ---, deslocando o custo da transição da lógica crítica para a
camada de apresentação. Em outras palavras, sob \emph{API-first}, a migração
deixou de ameaçar o núcleo funcional e passou a ser, predominantemente, um
exercício de portabilidade de interface.

O mesmo resultado fornece uma evidência, no plano do código-fonte, para a tese de
Karin (2025). Em revisão sistemática, Karin posiciona as soluções baseadas em
\emph{WebView} --- incluindo PWAs --- como as de maior \emph{overhead}, atribuindo
o custo à camada intermediária de DOM e do navegador \cite{karin2025}. As 1.239
linhas de utilitários eliminadas na migração (\emph{polyfills}, detecção de
dispositivo móvel e otimizações de rolagem) existiam exclusivamente para compensar
limitações do \emph{WebView} em navegadores móveis e correspondiam a 49\% do código
do cliente web. Sua supressão integral na arquitetura de renderização nativa
materializa, no artefato, o \emph{overhead} que Karin descreve de forma agregada:
além do custo de execução (energia e CPU), o paradigma \emph{WebView} impunha um
\emph{custo de engenharia} --- quase metade da base de código do cliente dedicada a
contornos ---, aqui suprimido. O achado, assim, corrobora empiricamente, e em um
caso real de saúde pública, a conclusão da revisão de Karin.

Por fim, os achados de Sokolov (2025) orientam a leitura dos resultados da Frente~A. Em \emph{benchmark} controlado via ADB, Sokolov observou que a
implementação nativa superou o React Native em eficiência de CPU e, sobretudo, no
crescimento de memória ao longo do tempo em rotinas contínuas, com vantagens
pontuais do React Native em cenários específicos \cite{sokolov2025}. Espera-se,
portanto, que a vantagem do cliente nativo sobre o PWA seja mais nítida em métricas
sensíveis à execução sustentada --- notadamente a memória RAM sob interação
contínua (rolagem) --- e mais modesta em tarefas dominadas por I/O e rede, como a
submissão do PDF, cuja latência é governada pela API. Diferentemente de Sokolov,
cujo \emph{benchmark} é sintético (criptografia AES) e desvinculado de domínio, a
Frente~A avaliará essas métricas em um fluxo real de uso clínico; a confirmação ou
refutação desse padrão de consumo de RAM constituirá o contraste central entre o
esperado pela literatura e o obtido em campo. As primeiras medições já apontam
nessa direção: em tablet Android, o cliente nativo consumiu \textbf{63,9\% menos
memória} que o PWA (84,9~MB contra 235,0~MB; $n=10$), convergindo com a vantagem
de memória relatada por Sokolov; a rolagem do cliente nativo mostrou-se fluida (apenas 0,2\% de quadros
\emph{janky}), ao passo que a aferição equivalente no PWA esbarrou em uma limitação
do instrumento, já que o Chrome não expõe essas estatísticas ao \texttt{gfxinfo}.

\subsection{Validade, limitações e extensão}
\label{sub:validade}

A Frente~B baseia-se na inspeção direta do código-fonte, não sujeita a viés
amostral. As Frentes~A e C controlam variáveis por desenho (dispositivo único e
execução alternada em A; heurísticas consagradas e severidade padronizada em C).
Como limitações: (i)~a Frente~A foi medida em um único dispositivo --- um tablet
de gama intermediária (Galaxy Tab~S9~FE), que não representa o aparelho de entrada
típico do SUS ---, de modo que as métricas \emph{absolutas} não se generalizam a
outros perfis de \emph{hardware}, preservando-se, contudo, a validade do resultado
\emph{relativo} (PWA $\times$ nativo medidos no mesmo aparelho); (ii)~a Frente~B é uma reconstrução documental
retrospectiva do mapa de migração; e (iii)~a Frente~C capta problemas de
conformidade heurística por especialistas, não a aceitação subjetiva de usuários
finais, e foi conduzida como estudo-piloto com painel reduzido, o que limita a generalização dos achados. As conclusões relativas sobre
reuso e desempenho têm potencial de transferência para migrações arquiteturais
correlatas na \emph{stack} adotada.

%
% Este arquivo SUBSTITUI o antigo frente_b_engenharia_migracao.tex: a Frente B
% já está incluída aqui como subseção (5.3), com as mesmas tabelas reais.
% Chaves de citação usadas já existem no seu .bib: oliveira2023, rampinelli2026,
% nielsen1994, cardieri2018.
% =====================================================================

\section{Resultados Obtidos}
\label{sec:resultados}

Este capítulo apresenta os resultados da avaliação tripartite do artefato. Por
transparência metodológica, distinguem-se dois tipos de resultado: os
\textbf{já apurados} --- a Frente~B (engenharia da migração), derivada da análise
direta do código-fonte e, portanto, determinística e reprodutível --- e os
\textbf{pendentes de coleta empírica} --- a Frente~A (desempenho), que depende dos
ensaios com ADB em dispositivo físico, e a Frente~C (usabilidade), que depende das
inspeções dos especialistas. Para as frentes pendentes, apresentam-se o
instrumento de medição já construído, tabelas preparadas para os dados e uma
demonstração do método com dados sintéticos, explicitamente rotulados como
ilustrativos.

\subsection{O produto: aplicativo móvel}
\label{sub:produto}

O primeiro resultado é o próprio artefato: o cliente móvel em React Native/Expo,
que consome a mesma API FastAPI do PWA e preserva as funcionalidades da versão
web (entrada manual de valores, envio de laudo em PDF, verificação de
disponibilidade do servidor, exibição de resultados/\emph{briefing} e tabela de
comparação de referências). A caracterização técnica dessa migração constitui a
Frente~B (Seção~\ref{sub:frenteB}).

\subsection{Frente A --- Desempenho comparativo (ADB)}
\label{sub:frenteA}

A Frente~A compara PWA e cliente móvel sobre as mesmas tarefas (T1 inicialização;
T2 submissão de PDF; T3 rolagem; T4 renderização de tabelas), no mesmo
dispositivo Android, com dez execuções por tarefa e cliente, alternadas para
diluir efeitos térmicos e de processos concorrentes. As métricas --- tempo de
resposta, uso de CPU, memória (RAM), taxa de quadros (FPS) e energia estimada ---
são coletadas via Android Debug Bridge (ADB) e \texttt{dumpsys}, conforme o
protocolo de Oliveira et al. (2023) \cite{oliveira2023}.

Os ensaios foram conduzidos em um tablet Samsung Galaxy Tab~S9~FE
(modelo SM-X510), com Android~16, com dez repetições por medição.
A Tabela~\ref{tab:desempenho} consolida os resultados obtidos. O achado mais
expressivo diz respeito ao consumo de memória: o cliente móvel nativo utilizou,
em média, \textbf{84,9~MB}, ante \textbf{235,0~MB} do PWA executado no
navegador --- uma redução de \textbf{63,9\%}, isto é, o PWA consome cerca de
\textbf{2,8 vezes mais} memória para a mesma funcionalidade. A inicialização a
frio do aplicativo nativo situou-se em \textbf{241~ms} (desvio-padrão de
33~ms). A fluidez da rolagem foi aferida com o \texttt{dumpsys gfxinfo}: o cliente nativo
manteve a rolagem praticamente sem travamentos, com apenas \textbf{0,2\%} de
quadros com atraso (\emph{janky}) e tempo de renderização no 95º percentil em torno
de 9--12~ms, dentro do orçamento do \emph{display}. Esse resultado reproduziu-se em duas coletas independentes (0,2\% de quadros \emph{janky} em ambas; $n=20$ no total), reforçando sua robustez. A mesma medição não pôde ser
obtida para o PWA: o Chrome realiza a composição gráfica por um \emph{pipeline}
próprio, que não expõe essas estatísticas ao \texttt{gfxinfo} --- uma assimetria de
medição já prevista para a comparação PWA $\times$ nativo. O uso de CPU durante a
rolagem manteve-se baixo em ambos os clientes, porém o instrumento
(\texttt{dumpsys cpuinfo}) não produziu variação apreciável entre repetições, não
permitindo comparação conclusiva nesse ponto. O consumo energético estimado,
também previsto no protocolo, não foi aferido nesta etapa, permanecendo como
trabalho futuro.

\begin{table}[H]
\centering
\caption{Desempenho PWA $\times$ Móvel (Galaxy Tab~S9~FE, Android~16; $n=10$; média $\pm$ desvio-padrão).}
\label{tab:desempenho}
\footnotesize
\begin{tabular}{|p{6.2cm}|c|c|c|}
\hline
\textbf{Métrica (direção)} & \textbf{PWA} & \textbf{Móvel} & \textbf{$\Delta$} \\ \hline
Memória RAM --- app carregado (MB, $\downarrow$) & $235{,}0 \pm 18{,}7$ & $84{,}9 \pm 0{,}2$ & $-63{,}9\%$ \\ \hline
Inicialização a frio (ms, $\downarrow$) & --- & $241{,}3 \pm 33{,}3$ & --- \\ \hline
Fluidez em rolagem --- quadros \emph{janky} (\%, $\downarrow$) & n/d (ver texto) & $0{,}2$ & --- \\ \hline
CPU em rolagem (\%, $\downarrow$) & \multicolumn{3}{c|}{baixo em ambos; sem diferença conclusiva} \\ \hline
\end{tabular}
\end{table}

A Figura~\ref{fig:desempenho} apresenta os resultados por métrica (um painel por
métrica; barras = média, hastes = desvio-padrão). O contraste de memória entre o
PWA e o cliente nativo é expressivo e corrobora, com dados primários de um caso
real, a tese de que o paradigma \emph{WebView} impõe maior consumo de recursos
\cite{karin2025, sokolov2025}.

\begin{figure}[H]
\centering
\includegraphics[width=0.95\textwidth]{fig_frente_a_desempenho_adb.png}
\caption{Desempenho PWA $\times$ Móvel por métrica (Galaxy Tab~S9~FE, Android~16; $n=10$; barras = média, hastes = desvio-padrão).}
\label{fig:desempenho}
\end{figure}

\subsection{Frente B --- Engenharia da migração}
\label{sub:frenteB}

A Frente~B caracteriza o reuso arquitetural viabilizado pela estratégia
\emph{API-first}, a partir do mapeamento de migração das duas bases de código do
cliente (PWA em React/TypeScript e aplicativo em React Native/Expo) e do serviço
compartilhado (\emph{backend} FastAPI). Conforme o procedimento descrito na
Seção~4.3, cada artefato preexistente foi classificado segundo três padrões de
portabilidade: \textbf{(i)}~reescrita completa; \textbf{(ii)}~portabilidade
estrutural com adaptação de elementos visuais; e \textbf{(iii)}~substituição por
componente nativo equivalente. As métricas de linhas de código (LOC) foram
extraídas diretamente do repositório.

\subsubsection{Reuso do núcleo funcional via API-first}

O resultado central é o \textbf{reuso integral (100\%) de toda a lógica
clinicamente sensível}. O \emph{backend} FastAPI --- extração por OCR/\emph{parsing},
classificação pelas faixas da PNS, seleção de especialidades e geração do
\emph{briefing} --- foi reaproveitado \textbf{sem uma única linha reescrita}, sendo
consumido pelo novo cliente pelos mesmos três \emph{endpoints} REST
(Tabela~\ref{tab:endpoints}). A base de conhecimento (\texttt{data/*.csv}) também
permaneceu inalterada. Em números, \textbf{1.654 linhas de Python} e \textbf{145
linhas de regras em CSV} foram preservadas na íntegra \cite{rampinelli2026}.

\begin{table}[H]
\centering
\caption{\emph{Endpoints} REST reaproveitados sem alteração pelos dois clientes.}
\label{tab:endpoints}
\footnotesize
\begin{tabular}{|l|l|p{6.2cm}|}
\hline
\textbf{Endpoint} & \textbf{Método} & \textbf{Função (reusada por PWA e Mobile)} \\ \hline
\texttt{/health} & GET & Verificação de disponibilidade do serviço \\ \hline
\texttt{/interpret} & POST & Interpretação a partir de laudo em PDF (multipart) \\ \hline
\texttt{/interpret-manual} & POST & Interpretação a partir de valores digitados (JSON) \\ \hline
\end{tabular}
\end{table}

\subsubsection{Mapa de migração dos componentes}

A Tabela~\ref{tab:mapa_migracao} apresenta o mapa de migração, artefato a
artefato. A maioria dos componentes de interface enquadrou-se no padrão~(ii):
estrutura declarativa e lógica de estado preservadas, com adaptação apenas das
primitivas visuais (de \texttt{<div>}/\texttt{<input>} para
\texttt{<View>}/\texttt{<TextInput>}/\texttt{StyleSheet}). As substituições
nativas (iii) concentraram-se nos pontos de contato com \emph{hardware} e ambiente.

\begin{table}[H]
\centering
\caption{Mapa de migração PWA $\rightarrow$ Mobile por artefato (LOC = linhas de código).}
\label{tab:mapa_migracao}
\scriptsize
\setlength{\tabcolsep}{3pt}
\begin{tabular}{|p{4.3cm}|p{4.0cm}|c|p{3.0cm}|}
\hline
\textbf{Artefato na PWA (LOC)} & \textbf{Contraparte no Mobile} & \textbf{LOC} & \textbf{Padrão} \\ \hline
\emph{Backend} FastAPI + serviços (1.654) & idêntico (via REST) & 1.654 & Reuso integral \\ \hline
Base de conhecimento \texttt{data/*.csv} (145) & idêntica (via REST) & 145 & Reuso integral \\ \hline
\texttt{types/index.ts} (56) & \texttt{types.ts} (58) & 58 & (ii) Estrutural (contrato) \\ \hline
\texttt{services/api.ts} (213) & \texttt{api.ts} (71) & 71 & (ii) Estrutural c/ adaptação de transporte \\ \hline
\texttt{App.tsx} (323) & \texttt{App.tsx} (209) & 209 & (ii) Estrutural (orquestração/estado) \\ \hline
\texttt{ManualEntryForm.tsx} (90) & \texttt{ManualEntryForm.tsx} & --- & (ii) Estrutural \\ \hline
\texttt{ResultsDisplay.tsx} (196) & \texttt{Results.tsx} & --- & (ii) Estrutural \\ \hline
\texttt{ReferenceComparison.tsx} (82) & \texttt{ReferenceComparisonTable.tsx} & --- & (ii) Estrutural \\ \hline
\texttt{FileUpload.tsx} (56, \emph{react-dropzone}) & \texttt{PdfUpload.tsx} (\emph{expo-document-picker}) & --- & (iii) Substituição nativa \\ \hline
\texttt{LoadingSpinner.tsx} (79) & \texttt{ActivityIndicator} (nativo) & --- & (iii) Substituição nativa \\ \hline
\texttt{index.tsx} (66, \emph{ReactDOM}) & \texttt{index.ts} (8, Expo) & 8 & (iii) \emph{Bootstrap} de plataforma \\ \hline
\texttt{utils/} (1.239: contorno de \emph{webview}) & --- (não portados) & 0 & Eliminados \\ \hline
\end{tabular}
\end{table}

\subsubsection{Taxa de reuso e eliminação de código de contorno}

A Tabela~\ref{tab:taxa_reuso} sintetiza a taxa de reuso por camada. Destaca-se um
achado quantitativo relevante: a PWA mantinha três módulos utilitários
(\texttt{polyfills.ts}, \texttt{mobileOptimizations.ts} e
\texttt{mobileDetection.ts}), somando \textbf{1.239 linhas --- 49\% de todo o
código do cliente web (2.525 LOC)} ---, que existiam exclusivamente para mitigar
limitações do paradigma \emph{webview} em navegadores móveis. Na renderização
nativa do React Native, esses contornos tornaram-se desnecessários e foram
\textbf{integralmente eliminados}. Esse resultado é o correlato, no plano do
código, do \emph{overhead} de recursos que a literatura atribui às soluções
baseadas em \emph{webview} \cite{oliveira2023}.

\begin{table}[H]
\centering
\caption{Taxa de reuso por camada na migração PWA $\rightarrow$ Mobile.}
\label{tab:taxa_reuso}
\footnotesize
\begin{tabular}{|p{4.6cm}|c|p{5.2cm}|}
\hline
\textbf{Camada} & \textbf{Reuso} & \textbf{Observação} \\ \hline
Lógica de negócio (\emph{backend} + regras) & 100\% & 1.654 LOC Python + 145 CSV, sem alteração \\ \hline
Contrato de dados (interfaces TS) & $\approx$100\% & 5 interfaces de domínio migradas verbatim \\ \hline
Serviços de comunicação (rede) & 3/3 operações & Semântica preservada; transporte adaptado \\ \hline
Apresentação (componentes de UI) & Reescrita estrutural & Padrão (ii): estrutura/estado preservados \\ \hline
Utilitários de contorno de \emph{webview} & 0\% (eliminados) & 1.239 LOC descartadas \\ \hline
\end{tabular}
\end{table}

Confirma-se, assim, a hipótese de reuso: a separação de responsabilidades
\emph{API-first} preservou o núcleo funcional (ativo de maior valor e risco) e
localizou o esforço de reengenharia na camada de apresentação
\cite{rampinelli2026, oliveira2023}.

\subsection{Frente C --- Usabilidade (inspeção heurística)}
\label{sub:frenteC}

A Frente~C avaliou a usabilidade por inspeção heurística \cite{nielsen1994,
cardieri2018}: avaliadores inspecionaram, de forma independente, os dois clientes
(PWA e aplicativo) sobre os mesmos cenários, registrando violações às dez
heurísticas e classificando cada problema em uma escala de severidade de 0 a 4, no
modelo de \emph{um registro por problema}. Em razão do número de avaliadores
efetivamente reunido, esta frente é reportada como \textbf{estudo-piloto} de
caráter exploratório: participaram \textbf{quatro avaliadores}, dos quais três
registraram problemas de usabilidade (severidade $\geq 1$) e um registrou uma
observação classificada como severidade~0 (não contabilizada como problema);
registros de teste foram descartados.

A Tabela~\ref{tab:usabilidade} sintetiza os resultados, a
Tabela~\ref{tab:inventario_c} detalha os problemas e a
Figura~\ref{fig:usabilidade} mostra a distribuição por heurística.

\begin{table}[H]
\centering
\caption{Síntese comparativa da inspeção heurística (estudo-piloto; $n=3$ problemas).}
\label{tab:usabilidade}
\footnotesize
\begin{tabular}{|p{6.8cm}|c|c|}
\hline
\textbf{Métrica} & \textbf{PWA} & \textbf{Móvel} \\ \hline
Problemas registrados (severidade $\geq 1$) & 2 & 1 \\ \hline
Índice médio de severidade & 1,5 & 2,0 \\ \hline
Problemas de severidade 3--4 & 0 & 0 \\ \hline
Heurísticas violadas & H7, H8 & H7 \\ \hline
\end{tabular}
\end{table}

\begin{table}[H]
\centering
\caption{Inventário dos problemas de usabilidade identificados (estudo-piloto).}
\label{tab:inventario_c}
\scriptsize
\setlength{\tabcolsep}{3pt}
\begin{tabular}{|c|c|c|c|p{4.9cm}|p{4.0cm}|}
\hline
\textbf{\#} & \textbf{Cliente} & \textbf{Heur.} & \textbf{Sev.} & \textbf{Problema (resumo)} & \textbf{Recomendação} \\ \hline
1 & PWA & H8 & 1 & Tela de resultados poluída: \emph{tag} ``Normal'' redundante em todos os \emph{cards}; faltam unidades e valores de referência; dados de infraestrutura visíveis; sem divisão por séries. & Limpar os \emph{cards}; incluir unidades e referências; agrupar por séries. \\ \hline
2 & PWA & H7 & 2 & Achado que diverge apenas da referência clássica (não da PNS) não é explicado ao usuário. & Explicar o achado também nesse caso de divergência. \\ \hline
3 & Móvel & H7 & 2 & Após anexar o PDF, a aba de digitação manual permanece habilitada, permitindo entrada redundante. & Desabilitar a digitação manual quando há PDF anexado. \\ \hline
\end{tabular}
\end{table}

\begin{figure}[H]
\centering
\includegraphics[width=0.9\textwidth]{fig_frente_c_inspecao_heuristica.png}
\caption{Frente~C --- problemas por heurística, PWA $\times$ Móvel (estudo-piloto; $n=3$ problemas).}
\label{fig:usabilidade}
\end{figure}

Não se identificou nenhum problema de severidade alta (3--4): as ocorrências
situaram-se entre cosmética (1) e menor (2), sugerindo ausência de barreiras
graves de usabilidade nos fluxos inspecionados. Os achados concentraram-se na
\textbf{tela de apresentação de resultados} e em aspectos de \textbf{fluxo}:
poluição visual e ausência de informações essenciais de leitura, como unidades e
valores de referência --- heurística de estética e design minimalista (H8) ---, e
comportamentos pouco flexíveis, como a falta de explicação para achados que
divergem apenas da referência clássica e a permanência da digitação manual após o
anexo do PDF --- heurística de flexibilidade e eficiência de uso (H7). Registrou-se
ainda, no aplicativo, uma observação de severidade~0 (a tela de resultados abrir no
rodapé em vez do topo), não contabilizada como problema. Embora o índice médio de
severidade tenha sido ligeiramente maior no aplicativo (2,0 contra 1,5 do PWA), o
número reduzido de registros não autoriza comparação quantitativa conclusiva entre
os clientes; no formato piloto, o valor desta frente reside na natureza
\textbf{qualitativa e acionável} dos achados, que apontam melhorias concretas de
interface e pavimentam uma inspeção heurística mais ampla como trabalho futuro.

\subsection{Síntese: resultados obtidos $\times$ hipóteses}
\label{sub:sintese}

A Tabela~\ref{tab:sintese} relaciona os resultados às hipóteses do trabalho. A
hipótese de \textbf{reuso arquitetural} (Frente~B) está empiricamente confirmada;
as hipóteses de \textbf{desempenho} (A) e \textbf{usabilidade} (C) têm instrumentos
validados e aguardam a coleta.

\begin{table}[H]
\centering
\caption{Resultados obtidos frente às hipóteses de pesquisa.}
\label{tab:sintese}
\footnotesize
\begin{tabular}{|p{5.6cm}|p{3.0cm}|p{5.0cm}|}
\hline
\textbf{Dimensão (hipótese)} & \textbf{Situação} & \textbf{Evidência} \\ \hline
Artefato móvel funcional & Alcançado & Seção~\ref{sub:produto} \\ \hline
Ganho de desempenho (A) & Confirmado (parcial) & Seção~\ref{sub:frenteA}: RAM $-$63,9\%; rolagem nativa fluida (0,2\% \emph{janky}) \\ \hline
Reuso arquitetural API-first (B) & Alcançado (real) & Seção~\ref{sub:frenteB}: 100\% da lógica; $-$1.239 LOC de \emph{webview} \\ \hline
Ganho de usabilidade (C) & Explorado (estudo-piloto) & Seção~\ref{sub:frenteC}: 3 problemas (sev. 1--2); achados qualitativos \\ \hline
\end{tabular}
\end{table}

\subsection{Discussão: confronto com a literatura}
\label{sub:discussao}

Os resultados da Frente~B dialogam diretamente com os trabalhos relacionados
(Seção~3) e permitem confrontar o observado com o esperado. De início, confirma-se
a expectativa, derivada da literatura de arquitetura \emph{API-first}
\cite{rampinelli2026}, de que a centralização da lógica no \emph{backend}
minimizaria o custo de uma migração de paradigma: a hipótese de reuso não apenas se
sustentou, como se mostrou integral (100\%) no núcleo funcional.

Esse achado oferece uma contrapartida quantitativa ao estudo qualitativo de Zhou
(2024). Enquanto os vinte desenvolvedores entrevistados por Zhou relataram a
migração entre plataformas e a manutenção como desafios recorrentes do
desenvolvimento \emph{cross-platform}, sem, contudo, quantificar o esforço nem
documentar um caso real de transição \cite{zhou2024}, a Frente~B evidencia que esse
atrito é fortemente modulado pela decisão arquitetural. No InterpreteLabBR, a
arquitetura \emph{API-first} converteu o que a literatura descreve como
``dificuldade de migração'' em um esforço mensurável e localizado --- 100\% de
reuso da lógica de negócio, sem reescrita, e a eliminação de 1.239 linhas de
código de contorno ---, deslocando o custo da transição da lógica crítica para a
camada de apresentação. Em outras palavras, sob \emph{API-first}, a migração
deixou de ameaçar o núcleo funcional e passou a ser, predominantemente, um
exercício de portabilidade de interface.

O mesmo resultado fornece uma evidência, no plano do código-fonte, para a tese de
Karin (2025). Em revisão sistemática, Karin posiciona as soluções baseadas em
\emph{WebView} --- incluindo PWAs --- como as de maior \emph{overhead}, atribuindo
o custo à camada intermediária de DOM e do navegador \cite{karin2025}. As 1.239
linhas de utilitários eliminadas na migração (\emph{polyfills}, detecção de
dispositivo móvel e otimizações de rolagem) existiam exclusivamente para compensar
limitações do \emph{WebView} em navegadores móveis e correspondiam a 49\% do código
do cliente web. Sua supressão integral na arquitetura de renderização nativa
materializa, no artefato, o \emph{overhead} que Karin descreve de forma agregada:
além do custo de execução (energia e CPU), o paradigma \emph{WebView} impunha um
\emph{custo de engenharia} --- quase metade da base de código do cliente dedicada a
contornos ---, aqui suprimido. O achado, assim, corrobora empiricamente, e em um
caso real de saúde pública, a conclusão da revisão de Karin.

Por fim, os achados de Sokolov (2025) orientam a leitura dos resultados da Frente~A. Em \emph{benchmark} controlado via ADB, Sokolov observou que a
implementação nativa superou o React Native em eficiência de CPU e, sobretudo, no
crescimento de memória ao longo do tempo em rotinas contínuas, com vantagens
pontuais do React Native em cenários específicos \cite{sokolov2025}. Espera-se,
portanto, que a vantagem do cliente nativo sobre o PWA seja mais nítida em métricas
sensíveis à execução sustentada --- notadamente a memória RAM sob interação
contínua (rolagem) --- e mais modesta em tarefas dominadas por I/O e rede, como a
submissão do PDF, cuja latência é governada pela API. Diferentemente de Sokolov,
cujo \emph{benchmark} é sintético (criptografia AES) e desvinculado de domínio, a
Frente~A avaliará essas métricas em um fluxo real de uso clínico; a confirmação ou
refutação desse padrão de consumo de RAM constituirá o contraste central entre o
esperado pela literatura e o obtido em campo. As primeiras medições já apontam
nessa direção: em tablet Android, o cliente nativo consumiu \textbf{63,9\% menos
memória} que o PWA (84,9~MB contra 235,0~MB; $n=10$), convergindo com a vantagem
de memória relatada por Sokolov; a rolagem do cliente nativo mostrou-se fluida (apenas 0,2\% de quadros
\emph{janky}), ao passo que a aferição equivalente no PWA esbarrou em uma limitação
do instrumento, já que o Chrome não expõe essas estatísticas ao \texttt{gfxinfo}.

\subsection{Validade, limitações e extensão}
\label{sub:validade}

A Frente~B baseia-se na inspeção direta do código-fonte, não sujeita a viés
amostral. As Frentes~A e C controlam variáveis por desenho (dispositivo único e
execução alternada em A; heurísticas consagradas e severidade padronizada em C).
Como limitações: (i)~a Frente~A foi medida em um único dispositivo --- um tablet
de gama intermediária (Galaxy Tab~S9~FE), que não representa o aparelho de entrada
típico do SUS ---, de modo que as métricas \emph{absolutas} não se generalizam a
outros perfis de \emph{hardware}, preservando-se, contudo, a validade do resultado
\emph{relativo} (PWA $\times$ nativo medidos no mesmo aparelho); (ii)~a Frente~B é uma reconstrução documental
retrospectiva do mapa de migração; e (iii)~a Frente~C capta problemas de
conformidade heurística por especialistas, não a aceitação subjetiva de usuários
finais, e foi conduzida como estudo-piloto com painel reduzido, o que limita a generalização dos achados. As conclusões relativas sobre
reuso e desempenho têm potencial de transferência para migrações arquiteturais
correlatas na \emph{stack} adotada.

%
% Este arquivo SUBSTITUI o antigo frente_b_engenharia_migracao.tex: a Frente B
% já está incluída aqui como subseção (5.3), com as mesmas tabelas reais.
% Chaves de citação usadas já existem no seu .bib: oliveira2023, rampinelli2026,
% nielsen1994, cardieri2018.
% =====================================================================

\section{Resultados Obtidos}
\label{sec:resultados}

Este capítulo apresenta os resultados da avaliação tripartite do artefato. Por
transparência metodológica, distinguem-se dois tipos de resultado: os
\textbf{já apurados} --- a Frente~B (engenharia da migração), derivada da análise
direta do código-fonte e, portanto, determinística e reprodutível --- e os
\textbf{pendentes de coleta empírica} --- a Frente~A (desempenho), que depende dos
ensaios com ADB em dispositivo físico, e a Frente~C (usabilidade), que depende das
inspeções dos especialistas. Para as frentes pendentes, apresentam-se o
instrumento de medição já construído, tabelas preparadas para os dados e uma
demonstração do método com dados sintéticos, explicitamente rotulados como
ilustrativos.

\subsection{O produto: aplicativo móvel}
\label{sub:produto}

O primeiro resultado é o próprio artefato: o cliente móvel em React Native/Expo,
que consome a mesma API FastAPI do PWA e preserva as funcionalidades da versão
web (entrada manual de valores, envio de laudo em PDF, verificação de
disponibilidade do servidor, exibição de resultados/\emph{briefing} e tabela de
comparação de referências). A caracterização técnica dessa migração constitui a
Frente~B (Seção~\ref{sub:frenteB}).

\subsection{Frente A --- Desempenho comparativo (ADB)}
\label{sub:frenteA}

A Frente~A compara PWA e cliente móvel sobre as mesmas tarefas (T1 inicialização;
T2 submissão de PDF; T3 rolagem; T4 renderização de tabelas), no mesmo
dispositivo Android, com dez execuções por tarefa e cliente, alternadas para
diluir efeitos térmicos e de processos concorrentes. As métricas --- tempo de
resposta, uso de CPU, memória (RAM), taxa de quadros (FPS) e energia estimada ---
são coletadas via Android Debug Bridge (ADB) e \texttt{dumpsys}, conforme o
protocolo de Oliveira et al. (2023) \cite{oliveira2023}.

Os ensaios foram conduzidos em um tablet Samsung Galaxy Tab~S9~FE
(modelo SM-X510), com Android~16, com dez repetições por medição.
A Tabela~\ref{tab:desempenho} consolida os resultados obtidos. O achado mais
expressivo diz respeito ao consumo de memória: o cliente móvel nativo utilizou,
em média, \textbf{84,9~MB}, ante \textbf{235,0~MB} do PWA executado no
navegador --- uma redução de \textbf{63,9\%}, isto é, o PWA consome cerca de
\textbf{2,8 vezes mais} memória para a mesma funcionalidade. A inicialização a
frio do aplicativo nativo situou-se em \textbf{241~ms} (desvio-padrão de
33~ms). A fluidez da rolagem foi aferida com o \texttt{dumpsys gfxinfo}: o cliente nativo
manteve a rolagem praticamente sem travamentos, com apenas \textbf{0,2\%} de
quadros com atraso (\emph{janky}) e tempo de renderização no 95º percentil em torno
de 9--12~ms, dentro do orçamento do \emph{display}. Esse resultado reproduziu-se em duas coletas independentes (0,2\% de quadros \emph{janky} em ambas; $n=20$ no total), reforçando sua robustez. A mesma medição não pôde ser
obtida para o PWA: o Chrome realiza a composição gráfica por um \emph{pipeline}
próprio, que não expõe essas estatísticas ao \texttt{gfxinfo} --- uma assimetria de
medição já prevista para a comparação PWA $\times$ nativo. O uso de CPU durante a
rolagem manteve-se baixo em ambos os clientes, porém o instrumento
(\texttt{dumpsys cpuinfo}) não produziu variação apreciável entre repetições, não
permitindo comparação conclusiva nesse ponto. O consumo energético estimado,
também previsto no protocolo, não foi aferido nesta etapa, permanecendo como
trabalho futuro.

\begin{table}[H]
\centering
\caption{Desempenho PWA $\times$ Móvel (Galaxy Tab~S9~FE, Android~16; $n=10$; média $\pm$ desvio-padrão).}
\label{tab:desempenho}
\footnotesize
\begin{tabular}{|p{6.2cm}|c|c|c|}
\hline
\textbf{Métrica (direção)} & \textbf{PWA} & \textbf{Móvel} & \textbf{$\Delta$} \\ \hline
Memória RAM --- app carregado (MB, $\downarrow$) & $235{,}0 \pm 18{,}7$ & $84{,}9 \pm 0{,}2$ & $-63{,}9\%$ \\ \hline
Inicialização a frio (ms, $\downarrow$) & --- & $241{,}3 \pm 33{,}3$ & --- \\ \hline
Fluidez em rolagem --- quadros \emph{janky} (\%, $\downarrow$) & n/d (ver texto) & $0{,}2$ & --- \\ \hline
CPU em rolagem (\%, $\downarrow$) & \multicolumn{3}{c|}{baixo em ambos; sem diferença conclusiva} \\ \hline
\end{tabular}
\end{table}

A Figura~\ref{fig:desempenho} apresenta os resultados por métrica (um painel por
métrica; barras = média, hastes = desvio-padrão). O contraste de memória entre o
PWA e o cliente nativo é expressivo e corrobora, com dados primários de um caso
real, a tese de que o paradigma \emph{WebView} impõe maior consumo de recursos
\cite{karin2025, sokolov2025}.

\begin{figure}[H]
\centering
\includegraphics[width=0.95\textwidth]{fig_frente_a_desempenho_adb.png}
\caption{Desempenho PWA $\times$ Móvel por métrica (Galaxy Tab~S9~FE, Android~16; $n=10$; barras = média, hastes = desvio-padrão).}
\label{fig:desempenho}
\end{figure}

\subsection{Frente B --- Engenharia da migração}
\label{sub:frenteB}

A Frente~B caracteriza o reuso arquitetural viabilizado pela estratégia
\emph{API-first}, a partir do mapeamento de migração das duas bases de código do
cliente (PWA em React/TypeScript e aplicativo em React Native/Expo) e do serviço
compartilhado (\emph{backend} FastAPI). Conforme o procedimento descrito na
Seção~4.3, cada artefato preexistente foi classificado segundo três padrões de
portabilidade: \textbf{(i)}~reescrita completa; \textbf{(ii)}~portabilidade
estrutural com adaptação de elementos visuais; e \textbf{(iii)}~substituição por
componente nativo equivalente. As métricas de linhas de código (LOC) foram
extraídas diretamente do repositório.

\subsubsection{Reuso do núcleo funcional via API-first}

O resultado central é o \textbf{reuso integral (100\%) de toda a lógica
clinicamente sensível}. O \emph{backend} FastAPI --- extração por OCR/\emph{parsing},
classificação pelas faixas da PNS, seleção de especialidades e geração do
\emph{briefing} --- foi reaproveitado \textbf{sem uma única linha reescrita}, sendo
consumido pelo novo cliente pelos mesmos três \emph{endpoints} REST
(Tabela~\ref{tab:endpoints}). A base de conhecimento (\texttt{data/*.csv}) também
permaneceu inalterada. Em números, \textbf{1.654 linhas de Python} e \textbf{145
linhas de regras em CSV} foram preservadas na íntegra \cite{rampinelli2026}.

\begin{table}[H]
\centering
\caption{\emph{Endpoints} REST reaproveitados sem alteração pelos dois clientes.}
\label{tab:endpoints}
\footnotesize
\begin{tabular}{|l|l|p{6.2cm}|}
\hline
\textbf{Endpoint} & \textbf{Método} & \textbf{Função (reusada por PWA e Mobile)} \\ \hline
\texttt{/health} & GET & Verificação de disponibilidade do serviço \\ \hline
\texttt{/interpret} & POST & Interpretação a partir de laudo em PDF (multipart) \\ \hline
\texttt{/interpret-manual} & POST & Interpretação a partir de valores digitados (JSON) \\ \hline
\end{tabular}
\end{table}

\subsubsection{Mapa de migração dos componentes}

A Tabela~\ref{tab:mapa_migracao} apresenta o mapa de migração, artefato a
artefato. A maioria dos componentes de interface enquadrou-se no padrão~(ii):
estrutura declarativa e lógica de estado preservadas, com adaptação apenas das
primitivas visuais (de \texttt{<div>}/\texttt{<input>} para
\texttt{<View>}/\texttt{<TextInput>}/\texttt{StyleSheet}). As substituições
nativas (iii) concentraram-se nos pontos de contato com \emph{hardware} e ambiente.

\begin{table}[H]
\centering
\caption{Mapa de migração PWA $\rightarrow$ Mobile por artefato (LOC = linhas de código).}
\label{tab:mapa_migracao}
\scriptsize
\setlength{\tabcolsep}{3pt}
\begin{tabular}{|p{4.3cm}|p{4.0cm}|c|p{3.0cm}|}
\hline
\textbf{Artefato na PWA (LOC)} & \textbf{Contraparte no Mobile} & \textbf{LOC} & \textbf{Padrão} \\ \hline
\emph{Backend} FastAPI + serviços (1.654) & idêntico (via REST) & 1.654 & Reuso integral \\ \hline
Base de conhecimento \texttt{data/*.csv} (145) & idêntica (via REST) & 145 & Reuso integral \\ \hline
\texttt{types/index.ts} (56) & \texttt{types.ts} (58) & 58 & (ii) Estrutural (contrato) \\ \hline
\texttt{services/api.ts} (213) & \texttt{api.ts} (71) & 71 & (ii) Estrutural c/ adaptação de transporte \\ \hline
\texttt{App.tsx} (323) & \texttt{App.tsx} (209) & 209 & (ii) Estrutural (orquestração/estado) \\ \hline
\texttt{ManualEntryForm.tsx} (90) & \texttt{ManualEntryForm.tsx} & --- & (ii) Estrutural \\ \hline
\texttt{ResultsDisplay.tsx} (196) & \texttt{Results.tsx} & --- & (ii) Estrutural \\ \hline
\texttt{ReferenceComparison.tsx} (82) & \texttt{ReferenceComparisonTable.tsx} & --- & (ii) Estrutural \\ \hline
\texttt{FileUpload.tsx} (56, \emph{react-dropzone}) & \texttt{PdfUpload.tsx} (\emph{expo-document-picker}) & --- & (iii) Substituição nativa \\ \hline
\texttt{LoadingSpinner.tsx} (79) & \texttt{ActivityIndicator} (nativo) & --- & (iii) Substituição nativa \\ \hline
\texttt{index.tsx} (66, \emph{ReactDOM}) & \texttt{index.ts} (8, Expo) & 8 & (iii) \emph{Bootstrap} de plataforma \\ \hline
\texttt{utils/} (1.239: contorno de \emph{webview}) & --- (não portados) & 0 & Eliminados \\ \hline
\end{tabular}
\end{table}

\subsubsection{Taxa de reuso e eliminação de código de contorno}

A Tabela~\ref{tab:taxa_reuso} sintetiza a taxa de reuso por camada. Destaca-se um
achado quantitativo relevante: a PWA mantinha três módulos utilitários
(\texttt{polyfills.ts}, \texttt{mobileOptimizations.ts} e
\texttt{mobileDetection.ts}), somando \textbf{1.239 linhas --- 49\% de todo o
código do cliente web (2.525 LOC)} ---, que existiam exclusivamente para mitigar
limitações do paradigma \emph{webview} em navegadores móveis. Na renderização
nativa do React Native, esses contornos tornaram-se desnecessários e foram
\textbf{integralmente eliminados}. Esse resultado é o correlato, no plano do
código, do \emph{overhead} de recursos que a literatura atribui às soluções
baseadas em \emph{webview} \cite{oliveira2023}.

\begin{table}[H]
\centering
\caption{Taxa de reuso por camada na migração PWA $\rightarrow$ Mobile.}
\label{tab:taxa_reuso}
\footnotesize
\begin{tabular}{|p{4.6cm}|c|p{5.2cm}|}
\hline
\textbf{Camada} & \textbf{Reuso} & \textbf{Observação} \\ \hline
Lógica de negócio (\emph{backend} + regras) & 100\% & 1.654 LOC Python + 145 CSV, sem alteração \\ \hline
Contrato de dados (interfaces TS) & $\approx$100\% & 5 interfaces de domínio migradas verbatim \\ \hline
Serviços de comunicação (rede) & 3/3 operações & Semântica preservada; transporte adaptado \\ \hline
Apresentação (componentes de UI) & Reescrita estrutural & Padrão (ii): estrutura/estado preservados \\ \hline
Utilitários de contorno de \emph{webview} & 0\% (eliminados) & 1.239 LOC descartadas \\ \hline
\end{tabular}
\end{table}

Confirma-se, assim, a hipótese de reuso: a separação de responsabilidades
\emph{API-first} preservou o núcleo funcional (ativo de maior valor e risco) e
localizou o esforço de reengenharia na camada de apresentação
\cite{rampinelli2026, oliveira2023}.

\subsection{Frente C --- Usabilidade (inspeção heurística)}
\label{sub:frenteC}

A Frente~C avaliou a usabilidade por inspeção heurística \cite{nielsen1994,
cardieri2018}: avaliadores inspecionaram, de forma independente, os dois clientes
(PWA e aplicativo) sobre os mesmos cenários, registrando violações às dez
heurísticas e classificando cada problema em uma escala de severidade de 0 a 4, no
modelo de \emph{um registro por problema}. Em razão do número de avaliadores
efetivamente reunido, esta frente é reportada como \textbf{estudo-piloto} de
caráter exploratório: participaram \textbf{quatro avaliadores}, dos quais três
registraram problemas de usabilidade (severidade $\geq 1$) e um registrou uma
observação classificada como severidade~0 (não contabilizada como problema);
registros de teste foram descartados.

A Tabela~\ref{tab:usabilidade} sintetiza os resultados, a
Tabela~\ref{tab:inventario_c} detalha os problemas e a
Figura~\ref{fig:usabilidade} mostra a distribuição por heurística.

\begin{table}[H]
\centering
\caption{Síntese comparativa da inspeção heurística (estudo-piloto; $n=3$ problemas).}
\label{tab:usabilidade}
\footnotesize
\begin{tabular}{|p{6.8cm}|c|c|}
\hline
\textbf{Métrica} & \textbf{PWA} & \textbf{Móvel} \\ \hline
Problemas registrados (severidade $\geq 1$) & 2 & 1 \\ \hline
Índice médio de severidade & 1,5 & 2,0 \\ \hline
Problemas de severidade 3--4 & 0 & 0 \\ \hline
Heurísticas violadas & H7, H8 & H7 \\ \hline
\end{tabular}
\end{table}

\begin{table}[H]
\centering
\caption{Inventário dos problemas de usabilidade identificados (estudo-piloto).}
\label{tab:inventario_c}
\scriptsize
\setlength{\tabcolsep}{3pt}
\begin{tabular}{|c|c|c|c|p{4.9cm}|p{4.0cm}|}
\hline
\textbf{\#} & \textbf{Cliente} & \textbf{Heur.} & \textbf{Sev.} & \textbf{Problema (resumo)} & \textbf{Recomendação} \\ \hline
1 & PWA & H8 & 1 & Tela de resultados poluída: \emph{tag} ``Normal'' redundante em todos os \emph{cards}; faltam unidades e valores de referência; dados de infraestrutura visíveis; sem divisão por séries. & Limpar os \emph{cards}; incluir unidades e referências; agrupar por séries. \\ \hline
2 & PWA & H7 & 2 & Achado que diverge apenas da referência clássica (não da PNS) não é explicado ao usuário. & Explicar o achado também nesse caso de divergência. \\ \hline
3 & Móvel & H7 & 2 & Após anexar o PDF, a aba de digitação manual permanece habilitada, permitindo entrada redundante. & Desabilitar a digitação manual quando há PDF anexado. \\ \hline
\end{tabular}
\end{table}

\begin{figure}[H]
\centering
\includegraphics[width=0.9\textwidth]{fig_frente_c_inspecao_heuristica.png}
\caption{Frente~C --- problemas por heurística, PWA $\times$ Móvel (estudo-piloto; $n=3$ problemas).}
\label{fig:usabilidade}
\end{figure}

Não se identificou nenhum problema de severidade alta (3--4): as ocorrências
situaram-se entre cosmética (1) e menor (2), sugerindo ausência de barreiras
graves de usabilidade nos fluxos inspecionados. Os achados concentraram-se na
\textbf{tela de apresentação de resultados} e em aspectos de \textbf{fluxo}:
poluição visual e ausência de informações essenciais de leitura, como unidades e
valores de referência --- heurística de estética e design minimalista (H8) ---, e
comportamentos pouco flexíveis, como a falta de explicação para achados que
divergem apenas da referência clássica e a permanência da digitação manual após o
anexo do PDF --- heurística de flexibilidade e eficiência de uso (H7). Registrou-se
ainda, no aplicativo, uma observação de severidade~0 (a tela de resultados abrir no
rodapé em vez do topo), não contabilizada como problema. Embora o índice médio de
severidade tenha sido ligeiramente maior no aplicativo (2,0 contra 1,5 do PWA), o
número reduzido de registros não autoriza comparação quantitativa conclusiva entre
os clientes; no formato piloto, o valor desta frente reside na natureza
\textbf{qualitativa e acionável} dos achados, que apontam melhorias concretas de
interface e pavimentam uma inspeção heurística mais ampla como trabalho futuro.

\subsection{Síntese: resultados obtidos $\times$ hipóteses}
\label{sub:sintese}

A Tabela~\ref{tab:sintese} relaciona os resultados às hipóteses do trabalho. A
hipótese de \textbf{reuso arquitetural} (Frente~B) está empiricamente confirmada;
as hipóteses de \textbf{desempenho} (A) e \textbf{usabilidade} (C) têm instrumentos
validados e aguardam a coleta.

\begin{table}[H]
\centering
\caption{Resultados obtidos frente às hipóteses de pesquisa.}
\label{tab:sintese}
\footnotesize
\begin{tabular}{|p{5.6cm}|p{3.0cm}|p{5.0cm}|}
\hline
\textbf{Dimensão (hipótese)} & \textbf{Situação} & \textbf{Evidência} \\ \hline
Artefato móvel funcional & Alcançado & Seção~\ref{sub:produto} \\ \hline
Ganho de desempenho (A) & Confirmado (parcial) & Seção~\ref{sub:frenteA}: RAM $-$63,9\%; rolagem nativa fluida (0,2\% \emph{janky}) \\ \hline
Reuso arquitetural API-first (B) & Alcançado (real) & Seção~\ref{sub:frenteB}: 100\% da lógica; $-$1.239 LOC de \emph{webview} \\ \hline
Ganho de usabilidade (C) & Explorado (estudo-piloto) & Seção~\ref{sub:frenteC}: 3 problemas (sev. 1--2); achados qualitativos \\ \hline
\end{tabular}
\end{table}

\subsection{Discussão: confronto com a literatura}
\label{sub:discussao}

Os resultados da Frente~B dialogam diretamente com os trabalhos relacionados
(Seção~3) e permitem confrontar o observado com o esperado. De início, confirma-se
a expectativa, derivada da literatura de arquitetura \emph{API-first}
\cite{rampinelli2026}, de que a centralização da lógica no \emph{backend}
minimizaria o custo de uma migração de paradigma: a hipótese de reuso não apenas se
sustentou, como se mostrou integral (100\%) no núcleo funcional.

Esse achado oferece uma contrapartida quantitativa ao estudo qualitativo de Zhou
(2024). Enquanto os vinte desenvolvedores entrevistados por Zhou relataram a
migração entre plataformas e a manutenção como desafios recorrentes do
desenvolvimento \emph{cross-platform}, sem, contudo, quantificar o esforço nem
documentar um caso real de transição \cite{zhou2024}, a Frente~B evidencia que esse
atrito é fortemente modulado pela decisão arquitetural. No InterpreteLabBR, a
arquitetura \emph{API-first} converteu o que a literatura descreve como
``dificuldade de migração'' em um esforço mensurável e localizado --- 100\% de
reuso da lógica de negócio, sem reescrita, e a eliminação de 1.239 linhas de
código de contorno ---, deslocando o custo da transição da lógica crítica para a
camada de apresentação. Em outras palavras, sob \emph{API-first}, a migração
deixou de ameaçar o núcleo funcional e passou a ser, predominantemente, um
exercício de portabilidade de interface.

O mesmo resultado fornece uma evidência, no plano do código-fonte, para a tese de
Karin (2025). Em revisão sistemática, Karin posiciona as soluções baseadas em
\emph{WebView} --- incluindo PWAs --- como as de maior \emph{overhead}, atribuindo
o custo à camada intermediária de DOM e do navegador \cite{karin2025}. As 1.239
linhas de utilitários eliminadas na migração (\emph{polyfills}, detecção de
dispositivo móvel e otimizações de rolagem) existiam exclusivamente para compensar
limitações do \emph{WebView} em navegadores móveis e correspondiam a 49\% do código
do cliente web. Sua supressão integral na arquitetura de renderização nativa
materializa, no artefato, o \emph{overhead} que Karin descreve de forma agregada:
além do custo de execução (energia e CPU), o paradigma \emph{WebView} impunha um
\emph{custo de engenharia} --- quase metade da base de código do cliente dedicada a
contornos ---, aqui suprimido. O achado, assim, corrobora empiricamente, e em um
caso real de saúde pública, a conclusão da revisão de Karin.

Por fim, os achados de Sokolov (2025) orientam a leitura dos resultados da Frente~A. Em \emph{benchmark} controlado via ADB, Sokolov observou que a
implementação nativa superou o React Native em eficiência de CPU e, sobretudo, no
crescimento de memória ao longo do tempo em rotinas contínuas, com vantagens
pontuais do React Native em cenários específicos \cite{sokolov2025}. Espera-se,
portanto, que a vantagem do cliente nativo sobre o PWA seja mais nítida em métricas
sensíveis à execução sustentada --- notadamente a memória RAM sob interação
contínua (rolagem) --- e mais modesta em tarefas dominadas por I/O e rede, como a
submissão do PDF, cuja latência é governada pela API. Diferentemente de Sokolov,
cujo \emph{benchmark} é sintético (criptografia AES) e desvinculado de domínio, a
Frente~A avaliará essas métricas em um fluxo real de uso clínico; a confirmação ou
refutação desse padrão de consumo de RAM constituirá o contraste central entre o
esperado pela literatura e o obtido em campo. As primeiras medições já apontam
nessa direção: em tablet Android, o cliente nativo consumiu \textbf{63,9\% menos
memória} que o PWA (84,9~MB contra 235,0~MB; $n=10$), convergindo com a vantagem
de memória relatada por Sokolov; a rolagem do cliente nativo mostrou-se fluida (apenas 0,2\% de quadros
\emph{janky}), ao passo que a aferição equivalente no PWA esbarrou em uma limitação
do instrumento, já que o Chrome não expõe essas estatísticas ao \texttt{gfxinfo}.

\subsection{Validade, limitações e extensão}
\label{sub:validade}

A Frente~B baseia-se na inspeção direta do código-fonte, não sujeita a viés
amostral. As Frentes~A e C controlam variáveis por desenho (dispositivo único e
execução alternada em A; heurísticas consagradas e severidade padronizada em C).
Como limitações: (i)~a Frente~A foi medida em um único dispositivo --- um tablet
de gama intermediária (Galaxy Tab~S9~FE), que não representa o aparelho de entrada
típico do SUS ---, de modo que as métricas \emph{absolutas} não se generalizam a
outros perfis de \emph{hardware}, preservando-se, contudo, a validade do resultado
\emph{relativo} (PWA $\times$ nativo medidos no mesmo aparelho); (ii)~a Frente~B é uma reconstrução documental
retrospectiva do mapa de migração; e (iii)~a Frente~C capta problemas de
conformidade heurística por especialistas, não a aceitação subjetiva de usuários
finais, e foi conduzida como estudo-piloto com painel reduzido, o que limita a generalização dos achados. As conclusões relativas sobre
reuso e desempenho têm potencial de transferência para migrações arquiteturais
correlatas na \emph{stack} adotada.
